% !PRINTED 411
\weberpage{411}
%<<CONT
chung $\Phi$, mit der alternirenden Permutationsgruppe von $n$ Ziffern
(ein- oder mehrstufig) isomorph. Die Gruppe der linearen Substitutionen
 $\chi(\vartheta)$ möge mit $L$ bezeichnet sein. Sie muss, da sie
endlich ist, mit einer der im siebenten Abschnitte betrachteten
Polyëdergruppen identisch sein.

Zur Vereinfachung machen wir nun von dem in \S.\,51, 2.
bewiesenen Satze Gebrauch, nach dem sich die Gruppe $L$ so
transformiren lässt, dass eine beliebige der nicht identischen
Substitutionen von $L$ eine Multiplication wird. Eine Transformation
 der Gruppe $L$ ist aber gleichbedeutend damit, dass für $\vartheta$
eine lineare Function von $\vartheta$ gesetzt wird, der dieselbe Eigenschaft
 wie der ursprünglichen Function $\vartheta$ zukommt.

Bezeichnen wir also mit $\vartheta_\pi$ die Function der Variablen
$x_0, x_1, \ldots, x_{n-1}$, die aus $\vartheta$ durch Anwendung der Permutation $\pi$
hervorgeht, so können wir $\vartheta$ so wählen, dass für eine bestimmte,
aber beliebige Permutation $\pi$
\[
\vartheta_\pi = \varepsilon\vartheta
\]
wird. Wendet man $\pi$ wiederholt an, so folgt
\[
\vartheta_{\pi^2} = \varepsilon^2\vartheta,\ \vartheta_{\pi^3} = \varepsilon^3\vartheta, \ldots,
\]
und wenn also $p$ der Grad der Permutation $\pi$ ist, so ist $\varepsilon$ eine
$p^{\text{te}}$ Einheitswurzel.

Ist zunächst $n = 3$, so besteht die alternirende Gruppe aus
den Potenzen der cyklischen Permutation $\gamma = (0, 1, 2)$. Wir
können annehmen, dass die der Permutation entsprechende Substitution
 $\chi(\vartheta)$ multiplicativ sei, dass also $\vartheta_\gamma = \varepsilon\vartheta$ sei, worin $\varepsilon$
eine dritte Einheitswurzel ist. $\vartheta^3$ ist daher eine alternirende
(oder symmetrische) Function. Als Resolvente mit einem Parameter
 erhalten wir also die reine Gleichung
\[
\vartheta^3 = z,
\]
wo $z$ eine alternirende Function ist. Wir können etwa für $\vartheta$
die \gesperrt{Lagrange}'sche Resolvente $x_0 + \varepsilon x_1 + \varepsilon^2 x_2$ nehmen, und
erhalten die bekannte Reduction der cubischen Gleichung auf
eine reine Gleichung (Bd. I, \S.\,159).

Wir gehen zu dem Falle $n = 4$ über, in dem die alternirende
 Gruppe $A$ die Permutationen
\[
1,\ \alpha_1 = (0, 1)(2, 3),\ \alpha_2 = (0, 2)(1, 3),\ \alpha_3 = (0, 3)(1, 2)
\]
enthält, die eine Vierergruppe $B$ bilden; ausserdem kommen
%CONT>>
